% -------------------------------------------------------
%  Section 4: Experimental setup
% -------------------------------------------------------
\section{بستر آزمایش و پیکربندی مقاله}
\label{sec:testbed}

آزمایش‌ها روی یک ماشین با پردازنده \lr{Xeon Platinum 8255C} و کارت گرافیک \lr{NVIDIA RTX 2080 Ti} با \num{۱۱} گیگابایت حافظه انجام شده‌اند. بهینه‌ساز در هر دو مدل گرادیان کاهشی تصادفی\pn{Stochastic Gradient Descent (SGD)} با تابع زیان آنتروپی متقاطع است و ابرپارامتر‌ها در جدول \ref{tab:hyper} آمده‌اند.

\begin{table}[!ht]
    \centering
    \footnotesize
    \caption{ابرپارامتر‌های گزارش‌شده مقاله (بازآوری جدول ۴ مقاله)}
    \label{tab:hyper}
    \begin{tabular}{|l|c|c|c|c|c|c|}
    \hline
    \textbf{مدل} & \textbf{دوره} & \textbf{اندازه دسته} & \textbf{نرخ یادگیری} & \textbf{زمان‌بند} & \textbf{کاهش وزن} & \textbf{تکانه} \\ \hline
    \lr{VGG16}    & ۲۰ & ۸  & \num{۰٫۰۰۱} & ندارد & \num{۰٫۰۰۰۵} & \num{۰٫۹} \\ \hline
    \lr{ResNet50} & ۱۵ & ۶۴ & \num{۰٫۰۱}  & نمایی با نرخ \num{۰٫۹} & \num{۰٫۰۰۶} & \num{۰٫۹} \\ \hline
    \end{tabular}
\end{table}

دو ملاحظه در همین جدول دیده می‌شود. نخست، کاهش وزن \num{۰٫۰۰۶} برای \lr{ResNet50} یک مرتبه بزرگی بیش از مقادیر متعارف در تنظیم دقیق شبکه‌های از پیش‌آموخته است و هیچ توضیحی درباره انتخاب آن داده نمی‌شود. دوم، تقریباً همه ابرپارامتر‌ها میان دو معماری متفاوت‌اند؛ بنابراین هر مقایسه‌ای میان \lr{VGG16} و \lr{ResNet50} در جدول نتایج، مقایسه دو \emph{رویه آموزشی} است و نه دو \emph{معماری}. مقاله تنها تفاوت اندازه دسته را با محدودیت حافظه توجیه می‌کند.

مقاله دو سنجه گزارش می‌کند. نخست صحت روی مجموعه آموزش، که خود نویسنده می‌پذیرد سنجه ضعیفی است. دوم زیان لگاریتمی چندکلاسه\pn{Multi-class Logarithmic Loss} روی مجموعه آزمون بدون برچسب که با ارسال پیش‌بینی‌ها به سرور کگل محاسبه می‌شود:
\begin{equation}
    \mathrm{logloss} = -\frac{1}{N}\sum_{i=1}^{N}\sum_{j=1}^{M} y_{ij}\log(p_{ij})
\end{equation}
که در آن $N$ تعداد نمونه‌های آزمون، $M$ تعداد خانواده‌ها، $y_{ij}$ نشانگر دودویی عضویت و $p_{ij}$ احتمال پیش‌بینی‌شده است. سرور کگل این مقدار را روی دو زیرمجموعه جدا می‌دهد: نمره عمومی\pn{Public Score} روی \pct{۳۰} نمونه‌ها و نمره خصوصی\pn{Private Score} روی \pct{۷۰} باقی‌مانده. مقاله نمره خصوصی را نتیجه نهایی می‌گیرد که انتخاب درستی است.

استفاده از یک سرور ارزیابی بیرونی با برچسب‌های پنهان، نقطه قوت روش‌شناختی مقاله است و امکان نشت داده آزمون را از میان می‌برد؛ ولی چون هر ۲۰ پیکربندی روی همان مجموعه آزمون ارزیابی و سپس بهترین‌ها گزارش شده‌اند، عدد گزینش‌شده به‌طور نظام‌مند خوش‌بینانه است. این نکته را در \S\ref{sec:critique} پی می‌گیریم.
